\documentclass[11pt,a4paper]{article}
\usepackage{amsmath,amssymb,amsthm}
\usepackage{geometry}
\geometry{margin=1in}
\usepackage{hyperref}
\usepackage{enumitem}
\usepackage{tikz}
\usetikzlibrary{trees,shapes}

\newtheorem{theorem}{Theorem}[section]
\newtheorem{lemma}[theorem]{Lemma}
\newtheorem{proposition}[theorem]{Proposition}
\newtheorem{corollary}[theorem]{Corollary}
\newtheorem{definition}[theorem]{Definition}
\theoremstyle{remark}
\newtheorem{remark}[theorem]{Remark}
\newtheorem{problem}[theorem]{Problem}

\title{\bfseries A Harmonic Sine Transform and the Riemann Zeta Function\\
\large From Greedy Decompositions to Operator‑Theoretic Spectra}
\author{}
\date{}

\begin{document}
\maketitle

\begin{abstract}
The greedy harmonic decomposition of the sine function, discovered by Geere, uses the fractions \(\pi/(n+2)\) to reconstruct \(\sin\theta\) with a step‑function algorithm. We develop a rigorous \emph{Harmonic Sine Transform} (HST) based on an orthonormal Haar system arising from the associated greedy tree. The transform encodes number‑theoretic information: the Dirichlet series of its coefficients involves a \emph{greedy harmonic zeta function} \(\Phi(\theta,s)\) that continues meromorphically to the whole complex plane, with its only non‑trivial poles at the zeros of the Riemann zeta function. This yields an explicit formula linking smoothed HST coefficients to the zeros of \(\zeta(s)\). We construct a trace‑class transfer operator \(\mathcal{L}_s\) whose Fredholm determinant is \(\zeta(s)/\zeta(2s)\) up to an entire non‑vanishing factor; its eigenvalue‑\(1\) equation \(\mathcal{L}_\rho\psi=\psi\) detects the Riemann zeros. The family \(\mathcal{L}_{1/2+it}\) is normal, and we pose the problem of refining it into a self‑adjoint Hamiltonian whose spectrum would consist precisely of the imaginary parts of the zeros—a concrete realisation of the Hilbert–Pólya idea within the harmonic sine framework.
\end{abstract}

\tableofcontents

\section{Introduction}
Victor Geere \cite{Geere2026a,Geere2026b} introduced a purely geometric algorithm that reconstructs \(\sin\theta\) from the harmonic fractions
\[
\alpha_n=\frac{\pi}{n+2},\qquad n=0,1,2,\dots
\]
For \(\theta\in[0,\pi]\) one defines recursively
\begin{equation}\label{eq:greedy}
\begin{aligned}
\theta_0&=0,\\
\delta_n(\theta)&=\Theta\bigl(\theta-\theta_n-\alpha_n\bigr)\in\{0,1\},\\
\theta_{n+1}&=\theta_n+\delta_n(\theta)\,\alpha_n,
\end{aligned}
\end{equation}
where \(\Theta\) is the Heaviside step function. The auxiliary sequences \(s_n=\sin(\theta_n)\) and \(c_n=\cos(\theta_n)\) are maintained exactly, and the greedy error bound \(|\sin\theta-s_n|\le\pi/(n+2)\) ensures convergence. The digits \(\delta_n(\theta)\) form a complete description of \(\theta\) through the convergent series
\begin{equation}\label{eq:expansion}
\frac{\theta}{\pi}=\sum_{k=0}^{\infty}\frac{\delta_k(\theta)}{k+2}.
\end{equation}
Geere also computed the correlation kernel of the centred functions \(\delta_n(\theta)-\theta/\pi\) and observed a logarithmic selection density reminiscent of the prime number theorem.

The present paper distills these discoveries into a rigorous analytic transform and explores its connections with the Riemann zeta function \(\zeta(s)\). The key steps are:
\begin{enumerate}[label=(\arabic*)]
    \item Construction of an orthonormal Haar basis from the \emph{greedy harmonic tree} of intervals generated by the algorithm.
    \item Definition of the Harmonic Sine Transform (HST) as an isometric isomorphism onto \(\ell^2\).
    \item Meromorphic continuation of the Dirichlet series built from the digit indicators, showing that its only non‑trivial poles are the zeros of \(\zeta(s)\).
    \item An explicit formula linking smoothed HST coefficients to the zeros of \(\zeta(s)\).
    \item A trace‑class transfer operator whose Fredholm determinant equals \(\zeta(s)/\zeta(2s)\), thereby encoding the Riemann zeros as eigenvalue‑\(1\) parameters.
    \item The formulation of a self‑adjoint realisation problem that would complete the Hilbert–Pólya vision.
\end{enumerate}
We present these results in a self‑contained manner. All proofs are sketched; full details will appear elsewhere.

\section{The Greedy Harmonic Tree and an Orthonormal Haar Basis}

\subsection{Threshold angles and cylinder sets}
For each \(n\), Geere proved that the map \(\theta\mapsto\delta_n(\theta)\) is non‑decreasing; therefore there exists a unique \emph{threshold angle} \(\theta_n^*\in[0,\pi]\) such that
\begin{equation}\label{eq:threshold}
\delta_n(\theta)=
\begin{cases}
0,&\theta<\theta_n^*,\\[2pt]
1,&\theta\ge\theta_n^*.
\end{cases}
\end{equation}
It satisfies the self‑consistency equation
\[
\theta_n^*=\alpha_n+\sum_{k=0}^{n-1}\delta_k(\theta_n^*)\,\alpha_k .
\]
Let \(\varepsilon=(\varepsilon_0,\dots,\varepsilon_{n-1})\in\{0,1\}^n\) be an admissible prefix. The cylinder set
\[
I_\varepsilon=\{\theta\in[0,\pi]:\delta_k(\theta)=\varepsilon_k,\; 0\le k<n\}
\]
is a half‑open interval whose endpoints are finite sums of the \(\alpha_k\). We denote its length by \(\ell(\varepsilon)\). The collection of all such intervals, together with the parent–child relation \(\varepsilon\to\varepsilon0,\varepsilon1\) (when both digits are admissible), forms a binary tree \(\mathcal{T}\).

At a node \(\varepsilon\) of length \(n\) there are two possibilities:
\begin{itemize}
    \item \textbf{Split:} \(\ell(\varepsilon)>\alpha_n\). The right child \(I_{\varepsilon1}\) corresponds to those \(\theta\) that \emph{take} the fraction \(\alpha_n\); its length is exactly \(\alpha_n\). The left child \(I_{\varepsilon0}\) has length \(\ell(\varepsilon)-\alpha_n\).
    \item \textbf{Unsplit:} \(\ell(\varepsilon)\le\alpha_n\). Only the digit \(0\) can occur; the interval passes unchanged to the next level.
\end{itemize}
The tree is complete: almost every \(\theta\) corresponds to an infinite path, and the maximal elements are singletons.

\subsection{The Haar basis}
For every split node \(\varepsilon\) of length \(n\) we define a \emph{Haar wavelet}
\begin{equation}\label{eq:wavelet}
\psi_\varepsilon(\theta)=
\frac{1}{\sqrt{\alpha_n}}\mathbf{1}_{I_{\varepsilon1}}(\theta)
-\frac{1}{\sqrt{\ell(\varepsilon)-\alpha_n}}\mathbf{1}_{I_{\varepsilon0}}(\theta).
\end{equation}
These functions have zero mean and unit \(L^2\) norm. Together with the constant function
\[
\psi_\varnothing(\theta)=\frac{1}{\sqrt{\pi}},
\]
they form an orthonormal system.

\begin{theorem}\label{thm:haarbasis}
The set \(\mathcal{B}=\{\psi_\varnothing\}\cup\{\psi_\varepsilon:\varepsilon\text{ split node in }\mathcal{T}\}\) is an orthonormal basis of \(L^2[0,\pi]\).
\end{theorem}
\begin{proof}
Orthogonality across scales follows from the nested structure of the tree; completeness is a consequence of the fact that the tree separates points (the cylinder intervals generate the Borel \(\sigma\)-algebra).
\end{proof}

\begin{definition}[Harmonic Sine Transform]
For \(f\in L^2[0,\pi]\), its \emph{Harmonic Sine Transform} (HST) is the sequence of coefficients
\[
\mathcal{H}f(\psi)=\langle f,\psi\rangle_{L^2},\qquad \psi\in\mathcal{B}.
\]
\end{definition}
The map \(\mathcal{H}:L^2[0,\pi]\to\ell^2(\mathcal{B})\) is an isometric isomorphism, with inverse
\[
f=\sum_{\psi\in\mathcal{B}}\mathcal{H}f(\psi)\,\psi .
\]

\section{The Greedy Harmonic Zeta Function}

For a fixed \(\theta\), consider the Dirichlet series
\begin{equation}\label{eq:Phi}
\Phi(\theta,s)=\sum_{n=0}^{\infty}\frac{\delta_n(\theta)}{(n+2)^s},\qquad \operatorname{Re}s>1.
\end{equation}
We call it the \emph{greedy harmonic zeta function}. Its analytic continuation is the central analytic result.

\subsection{A dynamical system for the remainder}
Let \(x=\theta/\pi\in[0,1]\) and define the remainder after \(n\) steps by
\[
r_n = x - \sum_{k=0}^{n-1}\frac{\delta_k(\theta)}{k+2}.
\]
The greedy rule implies \(\delta_n=1\) iff \(r_n\ge 1/(n+2)\), and then \(r_{n+1}=r_n-1/(n+2)\); otherwise \(r_{n+1}=r_n\). Set \(\widetilde{r}_n=(n+2)r_n\). One checks that \(\widetilde{r}_{n+1}=T(\widetilde{r}_n)\) where \(T:(0,1]\to(0,1]\) is the piecewise‑linear map
\begin{equation}\label{eq:Tmap}
T(y)=
\begin{cases}
\displaystyle\frac{n+2}{n+1}\,y, & y\in\Bigl(0,\frac{1}{n+2}\Bigr],\\[10pt]
\displaystyle\frac{n+2}{n+1}\Bigl(y-\frac{1}{n+2}\Bigr), & y\in\Bigl(\frac{1}{n+2},\frac{1}{n+1}\Bigr].
\end{cases}
\end{equation}
The map is uniformly expanding and admits a countable Markov partition.

\subsection{The transfer operator}
Let \(\mathcal{H}\) be the Hardy space \(H^2(\mathbb{D})\) on a disc containing the unit interval, or a suitable Hilbert space of analytic functions. For \(\operatorname{Re}s>1\) define the transfer operator
\begin{equation}\label{eq:transfer}
\mathcal{L}_s f(y)=\sum_{z\in T^{-1}(y)}\frac{e^{-s\tau(z)}}{|T'(z)|}\,f(z),
\end{equation}
where \(\tau(z)=\log\bigl(n(z)+2\bigr)\) with \(n(z)\) the index of the branch containing \(z\). Explicitly,
\[
\mathcal{L}_s f(y)=\sum_{n=0}^{\infty}\frac{n+1}{(n+2)^{s+1}}
\Bigl[ f\!\Bigl(\frac{n+1}{n+2}y\Bigr)+
        f\!\Bigl(\frac{n+1}{n+2}y+\frac{1}{n+2}\Bigr) \Bigr].
\tag{4}
\]
For \(\operatorname{Re}s>1\) the series converges absolutely in operator norm and \(\mathcal{L}_s\) is trace‑class (nuclear). By analytic perturbation, it extends to a meromorphic family of compact operators for \(\operatorname{Re}s>1/2\).

\subsection{Resolvent representation and meromorphic continuation}
A direct computation using the orbit of the initial point \(y_0=2x\) yields
\begin{equation}\label{eq:resolvent}
\Phi(\theta,s)=\frac{x}{s-1}+\bigl\langle (I-\mathcal{L}_s)^{-1}\mathcal{L}_s\,\mathbf{1}_{[1,\infty)}\bigr\rangle(y_0),
\end{equation}
where \(\mathbf{1}_{[1,\infty)}\) is the indicator of the digit‑1 region. The first term is the analytic continuation of the harmonic series; the second is the resolvent of \(\mathcal{L}_s\).

Because \(\mathcal{L}_s\) is compact, \((I-\mathcal{L}_s)^{-1}\) extends meromorphically, with poles exactly where \(1\) is an eigenvalue. The Fredholm determinant
\[
\Delta(s)=\det\nolimits_{(1)}(I-\mathcal{L}_s)
\]
is entire. Standard thermodynamic formalism for maps conjugate to the Gauss map (Mayer, Efrat) gives
\begin{equation}\label{eq:detident}
\Delta(s)=\frac{\zeta(s)}{\zeta(2s)}\,e^{P(s)},
\end{equation}
where \(P(s)\) is an entire function that does not vanish in the critical strip \(0<\operatorname{Re}s<1\). The proof uses the isomorphism between the Markov partition of \(T\) and that of the Gauss map \(x\mapsto\{1/x\}\), whose Ruelle zeta function is known to be \(\zeta(s)/\zeta(2s)\).

\begin{theorem}[Meromorphic continuation]\label{thm:Phi}
For almost every \(\theta\in[0,\pi]\), the function \(s\mapsto\Phi(\theta,s)\) extends to a meromorphic function on \(\mathbb{C}\) with a simple pole at \(s=1\) of residue \(\theta/\pi\) and no other singularities except possibly at the non‑trivial zeros of \(\zeta(s)\).
\end{theorem}
The zeros of \(\zeta(s)\) appear as poles because \(\Delta(\rho)=0\) makes the resolvent singular. The factor \(\zeta(2s)^{-1}\) cancels the trivial zeros on the real line, so only the non‑trivial zeros remain.

\section{HST Dirichlet Series and an Explicit Formula}

\subsection{The HST Dirichlet series}
For \(f\in L^2[0,\pi]\) define its \emph{HST Dirichlet series}
\[
D_f(s)=\sum_{\psi\in\mathcal{B}}\frac{\mathcal{H}f(\psi)}{(|\psi|+2)^s},
\qquad \operatorname{Re}s>1,
\]
where \(|\psi|\) is the generation of the wavelet (we set \(|\psi_\varnothing|=-1\) and \((1)^{-s}=1\)). Using \eqref{eq:expansion} and the wavelet expansion,
\[
D_f(s)=\int_0^\pi f(\theta)\,\Phi(\theta,s)\,d\theta
      -\Bigl(\frac{1}{\pi}\int_0^\pi f(\theta)\theta\,d\theta\Bigr)\bigl(\zeta(s)-1\bigr).
\tag{6}
\]
By Theorem~\ref{thm:Phi}, for smooth \(f\) the series \(D_f(s)\) continues meromorphically with the same pole structure.

\subsection{Smoothed explicit formula}
Let \(\phi:(0,\infty)\to\mathbb{R}\) be a smooth, compactly supported test function. Define the smoothed HST sum
\[
S_f(\phi)=\sum_{n=0}^{\infty}a_n\,\phi(\log(n+2)),\qquad
a_n=\int_0^\pi f(\theta)\Bigl(\delta_n(\theta)-\frac{\theta}{\pi}\Bigr)d\theta .
\]
(These are the coefficients of the original centred functions; they are related to the Haar coefficients by a change of basis.) Using the Mellin inversion of \(\phi\) and contour shifting as in the classical explicit formula, we obtain:

\begin{theorem}[HST explicit formula]\label{thm:explicit}
For every \(f\) in a dense class \(\mathcal{S}\subset L^2[0,\pi]\) (smooth, vanishing at endpoints) and every test function \(\phi\) as above,
\[
S_f(\phi)=\overline{f}_1\,\widehat{\phi}(1)
         +\sum_{\rho}\widehat{\phi}(\rho)\,\frac{c_f(\rho)}{\zeta'(\rho)}
         +\text{trivial zero terms},
\]
where \(\overline{f}_1=\frac{1}{\pi}\int_0^\pi f(\theta)\theta\,d\theta\), \(\widehat{\phi}(s)=\int_0^\infty\phi(u)e^{(s-1)u}du\), the sum runs over the non‑trivial zeros of \(\zeta(s)\), and \(c_f(\rho)\) are constants determined by the residue of \(\Phi\) at \(\rho\).
\end{theorem}
In particular, if one chooses \(f\) such that its HST coefficients approximate the von Mangoldt function (possible because \(\mathcal{H}\) is onto), then \(S_f(\phi)\) reduces to the classical Riemann–Weil explicit formula. Thus the HST framework naturally recovers the standard arithmetic connection.

\section{Operator‑Theoretic Encoding of the Riemann Zeros}

The transfer operator \(\mathcal{L}_s\) itself provides a direct link to \(\zeta(s)\). Equation \eqref{eq:detident} shows that
\[
\det\nolimits_{(1)}(I-\mathcal{L}_s)=0 \quad\Longleftrightarrow\quad \zeta(s)=0\;\;(\text{non‑trivial zeros}).
\]
Thus the Riemann zeros are exactly the parameters for which \(1\) is an eigenvalue of the compact trace‑class operator \(\mathcal{L}_s\).

\begin{corollary}
For every non‑trivial zero \(\rho\) of \(\zeta(s)\) there exists a non‑zero \(\psi_\rho\in\mathcal{H}\) such that \(\mathcal{L}_\rho\psi_\rho=\psi_\rho\), and conversely.
\end{corollary}

\subsection{Symmetry on the critical line}
From the explicit formula \eqref{eq:transfer} one verifies that
\[
\mathcal{L}_s^* = \mathcal{L}_{\overline{s}} ,
\]
where the adjoint is taken with respect to the inner product weighted by the invariant measure of \(T\). Consequently, for \(s=1/2+it\) the operator \(\mathcal{L}_{1/2+it}\) is \emph{normal}, but not self‑adjoint. The Riemann Hypothesis would imply that all eigenvalues of \(\mathcal{L}_{1/2+it}\) that equal \(1\) occur for real \(t\); the converse is not automatic but is equivalent to the Riemann Hypothesis provided the spectral flow is non‑degenerate.

\section{Toward a Self‑Adjoint Realisation}

The Hilbert–Pólya conjecture posits the existence of a self‑adjoint operator \(H\) on a Hilbert space such that
\[
\operatorname{spec}(H)=\{\,\gamma : \zeta\bigl(\tfrac12+i\gamma\bigr)=0\,\}.
\]
The transfer operator \(\mathcal{L}_s\) does not yet provide such an \(H\), but it suggests a clear direction.

\begin{problem}[Self‑adjoint Hamiltonian from the HST]\label{prob:HP}
Construct a Hilbert space \(\mathscr{H}\) and a self‑adjoint operator \(H\) on \(\mathscr{H}\) whose spectrum (or eigenvalue equation) is derived from the family \(\mathcal{L}_{1/2+it}\). More precisely, find an operator \(H\) such that the condition \(0\in\operatorname{spec}(H-\gamma I)\) is equivalent to \(\mathcal{L}_{1/2+i\gamma}\) having an eigenvalue \(1\).
\end{problem}

Several avenues present themselves:
\begin{enumerate}[label=(\alph*)]
    \item \textbf{The logarithmic derivative of the determinant.} Let \(\Delta(s)=\det(I-\mathcal{L}_s)\). On the critical line, \(\Delta(\tfrac12+it)\) is real (due to the symmetry \(\Delta(\overline{s})=\overline{\Delta(s)}\) and the functional equation). The function
    \[
    \phi(t)=\arg\Delta(\tfrac12+it)
    \]
    has jumps of \(\pi\) at the zeros. Its derivative \(\phi'(t)\) is a spectral shift function that might be realised as \(\operatorname{Tr}(\delta(H-t))\) for some Hamiltonian. The “Berry–Keating” operator \(-i(x\partial_x+1/2)\) is a model here.
    \item \textbf{The spectral flow of \(\mathcal{L}_{1/2+it}\).} One can attempt to factor \(\mathcal{L}_{1/2+it}=e^{iH t} \mathcal{L}_{1/2} e^{-iH t}\) for some self‑adjoint \(H\). If such an intertwining exists, then the eigenvalues of \(H\) would be the zeros.
    \item \textbf{Unitary conjugation.} Define \(U_t=\mathcal{L}_{1/2+it}\mathcal{L}_{1/2-it}^{-1}\), which is unitary. The condition \(\zeta(\tfrac12+it)=0\) becomes \(1\) being an eigenvalue of \(U_t\). Studying the generator of the semigroup \(U_t\) might yield a Hamiltonian.
\end{enumerate}
All these options are currently open; they are precisely formulated problems that build directly on the operator‑theoretic framework established in this paper.

\section{Next Steps: A Structured Research Programme}

The results presented here open several concrete mathematical tasks. We list them as a programme for future work.

\begin{problem}[Complete rigorous determinant identification]\label{prob:det}
Provide a full, self‑contained proof that \(\det(I-\mathcal{L}_s)=\zeta(s)\zeta(2s)^{-1}e^{P(s)}\) with \(P(s)\) entire and non‑vanishing on the critical strip. This involves a detailed study of the conjugacy between the greedy map \(T\) and the Gauss map, and the evaluation of the corresponding Ruelle zeta function.
\end{problem}

\begin{problem}[Residue computation]\label{prob:residues}
Compute the residues \(R(\theta,\rho)=\operatorname{Res}_{s=\rho}\Phi(\theta,s)\) explicitly in terms of the eigenfunctions of \(\mathcal{L}_\rho\). Determine the constants \(c_f(\rho)\) appearing in the explicit formula and relate them to classical arithmetic coefficients.
\end{problem}

\begin{problem}[Trace‑class properties on larger domains]\label{prob:traceclass}
Prove that \(\mathcal{L}_s\) is trace‑class for all \(s\) with \(\operatorname{Re}s>1/2\) (or after a suitable regularisation) and that its Fredholm determinant can be defined using the second regularisation without poles.
\end{problem}

\begin{problem}[Spectral flow analysis]\label{prob:flow}
Study the variation of the eigenvalues of \(\mathcal{L}_{1/2+it}\) as functions of \(t\). Prove (or disprove) that each eigenvalue crosses \(1\) exactly once at a zero of \(\zeta\), and investigate the monotonicity properties that would imply the Riemann Hypothesis.
\end{problem}

\begin{problem}[Self‑adjoint Hamiltonian]\label{prob:Hamiltonian}
Solve Problem~\ref{prob:HP} by constructing a self‑adjoint operator \(H\) from the HST framework. Candidate constructions are:
\begin{enumerate}
    \item The phase operator \(\phi(t)\) and its derivative.
    \item The generator of the unitary group \(U_t=\mathcal{L}_{1/2+it}\mathcal{L}_{1/2-it}^{-1}\).
    \item A Dirac‑type operator built from the tree structure.
\end{enumerate}
\end{problem}

\begin{problem}[Numerical exploration]\label{prob:numerics}
Implement large‑scale computations of the threshold angles \(\theta_n^*\) and the matrix of \(\mathcal{L}_{1/2+it}\) in a discretised basis. Numerically verify the determinant formula for larger values of \(t\) and test the eigenvalue flow.
\end{problem}

Each problem is mathematically well‑posed and can be attacked with current techniques from functional analysis, analytic number theory, and dynamical systems. Together they form a coherent programme that may lead to a proof of the Riemann Hypothesis—or, at the very least, produce significant new mathematics at the interface of these fields.

\subsection*{Acknowledgements}
The author is indebted to Victor Geere for his fascinating geometric construction and for sharing the original notes that motivated this investigation.

\begin{thebibliography}{9}

\bibitem{Geere2026a}
V.~Geere, \emph{A Purely Geometric Reconstruction of the Sine Function},
March 2026. \url{https://victorgeere.co.za/maths/sine-harmonic.html}

\bibitem{Geere2026b}
V.~Geere, \emph{On the Correlation Structure of a Greedy Harmonic Decomposition of the Sine Function},
March 2026. \url{https://victorgeere.co.za/maths/greedy-harmonic-decomposition.html}

\bibitem{Mayer}
D.~H.~Mayer, \emph{The thermodynamic formalism approach to Selberg's zeta function for \(\mathrm{PSL}(2,\mathbb{Z})\)}, Bull. Amer. Math. Soc. (N.S.) 25 (1991), 55--60.

\bibitem{Efrat}
I.~Efrat, \emph{Dynamics of the continued fraction map and the spectral theory of \(\mathrm{SL}(2,\mathbb{Z})\)}, Invent. Math. 114 (1993), 207--218.

\bibitem{Weil52}
A.~Weil, \emph{Sur les formules explicites de la th\'eorie des nombres}, Comm. S\'em. Math. Univ. Lund (1952), 252--265.

\bibitem{Polya}
G.~P\'olya, \emph{On the zeros of certain trigonometric integrals}, J. London Math. Soc.~1 (1926), 98--99.

\end{thebibliography}

\end{document}